% ============================================================
%  Números Primos — Apostila Premium | PratiqueJá
%  Engine: XeLaTeX
% ============================================================
\documentclass[12pt]{article}
\usepackage{../../pratiqueja}

% \hlm — destaque de resultado em modo-math (cor sem bold)
\newcommand{\hlm}[1]{{\color{darkblue}#1}}

% Células do crivo: primo / composto / unidade
\newcommand{\pc}[1]{\cellcolor{rulebg}\textcolor{darkblue}{\bfseries #1}}
\newcommand{\uc}[1]{\cellcolor{obsbg}\textcolor{mutedtext}{#1}}

\setsubject{Números Primos}

\begin{document}
\pagenumbering{arabic}
\setstretch{1.2}
\color{bodytext}

\heroheader{Números Primos}%
           {Crivo, Coprimos e Número de Divisores}%
           {Matemática · Básico}

\setlength{\columnsep}{18pt}
\setlength{\columnseprule}{0.5pt}
\def\columnseprulecolor{\color{exborder}}

\begin{multicols}{2}

%─────────────────────────────────────────────────────────────
\TW{Def}{badgepurple}{Primos e Compostos}{Divisores e classificação}

\regra{Um natural maior que $1$ é \textbf{primo} se seus únicos
divisores são \textbf{$1$ e ele mesmo}. Caso contrário é
\textbf{composto} (tem algum divisor além de $1$ e de si).}

\exemp{%
  \textbf{Primos:} $2, 3, 5, 7, 11, 13, 17, 19, 23, 29, 31\ldots$\\[3pt]
  \textbf{Compostos:} $4, 6, 8, 9, 10, 12, 14, 15\ldots$%
}

\nota{$1$ \textbf{não} é primo nem composto (convenção). E $2$ é o
\textbf{único} primo par.}

%─────────────────────────────────────────────────────────────
\TW{Crivo}{darkblue}{Crivo de Eratóstenes}{Encontrar todos os primos até $N$}

\regra{%
\textbf{1.} Liste de $2$ até $N$.\\[2pt]
\textbf{2.} Com $p=2$, marque como compostos os múltiplos de $p$
maiores que $p$.\\[2pt]
\textbf{3.} Avance para o próximo não marcado e repita.\\[2pt]
\textbf{4.} Pare quando $p^2 > N$. Os não marcados são primos.}

\SH{Primos até 30}
\begin{center}\vspace{1pt}
\setlength{\tabcolsep}{4pt}\renewcommand{\arraystretch}{1.25}
\small
\begin{tabular}{|*{10}{c|}}
\hline
\uc{1} & \pc{2} & \pc{3} & 4 & \pc{5} & 6 & \pc{7} & 8 & 9 & 10\\
\hline
\pc{11} & 12 & \pc{13} & 14 & 15 & 16 & \pc{17} & 18 & \pc{19} & 20\\
\hline
21 & 22 & \pc{23} & 24 & 25 & 26 & 27 & 28 & \pc{29} & 30\\
\hline
\end{tabular}
\end{center}

\exemp{Primos até 30: \hlm{2, 3, 5, 7, 11, 13, 17, 19, 23, 29}}

%─────────────────────────────────────────────────────────────
\TW{$\div p$}{darkblue}{Decomposição em Fatores Primos}{Todo composto é produto de primos}

\regra{\textbf{Teorema Fundamental da Aritmética:} todo composto se
escreve como produto de primos. Divida sempre pelo \textbf{menor}
primo que o divide, até chegar a $1$.}

\SH{Exemplo — decompor 60}
\begin{center}\vspace{1pt}
\renewcommand{\arraystretch}{1.2}\small
\begin{tabular}{r|l}
60 & 2\\
30 & 2\\
15 & 3\\
 5 & 5\\
\hline
 1 & \\
\end{tabular}
\end{center}
\exemp{$60 = 2^2 \times 3 \times 5$}

\SH{Exemplo — decompor 84}
\begin{center}\vspace{1pt}
\renewcommand{\arraystretch}{1.2}\small
\begin{tabular}{r|l}
84 & 2\\
42 & 2\\
21 & 3\\
 7 & 7\\
\hline
 1 & \\
\end{tabular}
\end{center}
\exemp{$84 = 2^2 \times 3 \times 7$}

%─────────────────────────────────────────────────────────────
\TW{?}{darkblue}{Teste de Primalidade}{Como verificar se um número é primo}

\regra{Para testar se $n$ é primo, verifique a divisibilidade pelos
primos $p \leq \sqrt{n}$. Se nenhum divide $n$, então $n$ é primo.}

\SH{Exemplo — verificar 97}
\exemp{%
  $\sqrt{97} \approx 9{,}8$ → testar primos até $9$: $2, 3, 5, 7$\\[2pt]
  $97$ não é divisível por nenhum deles.\\[2pt]
  Logo, \hlm{97 é primo}.%
}

\nota{\textbf{Por que só até $\sqrt{n}$?} Se $n = a\times b$ com
$a \leq b$, então $a \leq \sqrt{n}$. Basta procurar o fator menor.}

%─────────────────────────────────────────────────────────────
\TW{Co}{badgepurple}{Números Coprimos}{Sem fatores primos em comum}

\regra{Dois números são \textbf{coprimos} (primos entre si) quando
não compartilham nenhum fator primo — isto é, MDC $= 1$.}

\exemp{%
  $9 = 3^2$ e $16 = 2^4$ — sem fatores comuns\\
  MDC$(9,16) = \hlm{1}$ \ok\\[3pt]
  $35 = 5\times7$ e $48 = 2^4\times3$ — sem fatores comuns\\
  MDC$(35,48) = \hlm{1}$ \ok%
}

\nota{Números \textbf{consecutivos} são sempre coprimos:\\[2pt]
MDC$(13,14) = \hlm{1}$, pois $13$ é primo e não divide $14$.}

%─────────────────────────────────────────────────────────────
\TW{$\tau$}{darkblue}{Número de Divisores}{Fórmula a partir dos fatores primos}

\regra{Se $n = p_1^{a_1}\times p_2^{a_2}\times\cdots\times p_k^{a_k}$,
o número total de divisores é:}

\formula{$\tau(n) = (a_1{+}1)(a_2{+}1)\cdots(a_k{+}1)$}

Cada expoente $a_i$ contribui com $(a_i+1)$ escolhas: $0,1,\ldots,a_i$.

\SH{Exemplo — divisores de 60}
\exemp{%
  $60 = 2^2\times3^1\times5^1$\\[2pt]
  $\tau(60) = (2{+}1)(1{+}1)(1{+}1) = 3\times2\times2 = \hlm{12}$\\[2pt]
  $1,2,3,4,5,6,10,12,15,20,30,60$ \ok%
}

\SH{Exemplo — divisores de 72}
\exemp{%
  $72 = 2^3\times3^2$\\[2pt]
  $\tau(72) = (3{+}1)(2{+}1) = 4\times3 = \hlm{12}$%
}

\end{multicols}

\end{document}
